%&pdftex 
%%%%%%%%%%%%%%%%%%%%%%%%%%%%%%%%%%%%%%%%%%%%%%%%%%%%%%%%%%%%%%%%%%%%%%%%%%%
%                                                                         %
%   ULTIMA MODIFICACION: 31 DE MARZO DE 1992, POR MARIO PEREZ RIERA.      %
%   ENVIADO AL MATHEMATISCHE NACHRICHTEN EL 20 DE JUNIO DE 1991.          %
%   ACEPTADO MARZO 1992 (VERSION CORREGIDA CON INDICACIONES DEL REFEREE)  %
%                                                                         %
%%%%%%%%%%%%%%%%%%%%%%%%%%%%%%%%%%%%%%%%%%%%%%%%%%%%%%%%%%%%%%%%%%%%%%%%%%%

\magnification = 1200
\font \peq = cmr8
\font \peqbf = cmbx8
\font \cmcsc = cmcsc10

\font\msymnormal=msbm10
\font\msympeq=msbm7
\font\msymmuypeq=msbm5
\newfam\numeros
\textfont\numeros=\msymnormal
\scriptfont\numeros=\msympeq
\scriptscriptfont\numeros=\msymmuypeq
\def\num{\fam\numeros\msymnormal}
\def\N{{\num N}}
\def\R{{\num R}}
\font\tipo = cmr12 at 10pt
\tipo

\def \proof {{\it Proof.} }
\def \qed {\vrule height 5 pt width 5 pt depth 0 pt}

\centerline{\bf WEIGHTED WEAK BEHAVIOUR OF FOURIER-JACOBI SERIES}

\bigskip

\centerline{Jos\'e J. Guadalupe
   \footnote{$^1$}{\peq Dpto. de Matem\'aticas. Universidad de Zaragoza. 50009
Zaragoza. Spain.}, Mario P\'erez $^1$ and Juan L. Varona
   \footnote{$^2$}{\peq Dpto. de Matem\'atica Aplicada. Colegio Universitario de
La Rioja. 26001 Logro\~no. Spain.}}
\footnote{}{\peq E-mail: JVARONA@CC.UNIZAR.ES}
\footnote{}{\peq 1991 Mathematics Subject Classification: 42C10.}
\footnote{}{\peq Research supported by DGICYT under grant PB89-0181-C02-02}
\footnote{}{\peq THIS PAPER WERE PUBLISHED IN:
Math.\ Nachrichten {\peqbf 158} (1992), 161--174.}


\bigskip

{\narrower\narrower\noindent{\bf Abstract. }{\sl Let $w(x) = (1-x)^\alpha
(1+x)^\beta$ be a Jacobi weight on the interval $[-1,1]$ and $1 < p < \infty$.
If either $\alpha > -1/2$ or $\beta > -1/2$ and $p$ is an endpoint of the
interval of mean convergence of the associated Fourier-Jacobi series, we show
that the partial sum operators $S_n$ are uniformly bounded from $L^{p,1}$ to
$L^{p,\infty}$, thus extending a previous result for the case that both
$\alpha, \beta > -1/2$. For $\alpha, \beta > -1/2$, we study the weak and
restricted weak $(p,p)$-type of the weighted operators $f \longrightarrow u
S_n(u^{-1} f)$, where $u$ is also a Jacobi weight.}\smallskip}

\bigskip

{\noindent \bf \S 1. Introduction and main results.}

\medskip

\noindent Let $w$ be a Jacobi weight on the interval $[-1,1]$, that is,
$$
   w(x)=(1-x)^{\alpha}(1+x)^{\beta}, \qquad \alpha, \beta > -1
$$
and let $1 < p < \infty$; $S_nf$ stands for the $n$-th partial sum of the
Fourier series associated to the Jacobi polynomials, orthonormal on $[-1,1]$
with respect to $w$. It is well known that $S_nf$ converges to $f$ for every $f
\in L^p(w)$ if and only if the partial sum operators $S_n$ are uniformly
bounded in $L^p(w)$, i.e., there exists a constant $C>0$ such that

$$
   \Vert S_nf \Vert _{L^p(w)} \leq C \Vert f \Vert _{L^p(w)}
   \quad \forall n \geq 0, \ \forall f \in L^p(w) \eqno (1)
$$
(throughout this paper, we will denote by $C$ a constant independent of $f$,
$n$, etc., but not necessarily the same at each occurrence). Furthermore, there
exists an open interval $(p_0, p_1)$ such that this boundedness holds if and
only if $p$ belongs to $(p_0, p_1)$ (see [6]). The assumption that either
$\alpha > -1/2$ or $\beta > -1/2$ is equivalent to $1 < p_0 < p_1 < \infty$.
More precisely, in this case (1) holds if and only if
$$
   {4(\alpha+1) \over 2\alpha+3} < p < {4(\alpha+1) \over 2\alpha+1}
$$
when $\alpha \geq \beta$ (and the analogous inequality with $\alpha$ replaced
by $\beta$ if $\beta \geq \alpha$).

In this paper we examine the behaviour of $S_n$ at the endpoints of the
interval of mean convergence. In order to do this we need some classical
definitions and notations. Given a measure $\mu$ and $1 \leq p < \infty$, the
space $L^p_*(\mu) = L^{p,\infty}(\mu)$ is defined to be the space of measurable
functions such that
$$
   \Vert f \Vert_{L^p_*(\mu)} = \sup_{y > 0} 
   y \left[ \mu( \{ x : \vert f(x) \vert > y \} ) \right]^{1/p} < \infty.
$$
An operator $T$ is of weak $(p,p)$-type if $T : L^p(\mu) \longrightarrow
L^p_*(\mu)$ is bounded. Now, let $f^*$ be the nonincreasing rearrangement of
$f$, given by $f^*(t) = \inf \{s : \lambda(s) \leq t\}$, where $\lambda$
denotes the distribution function of $f$. Then, the Lorentz space
$L^{p,r}(\mu)$ is the class of all measurable functions $f$ satisfying
$$
   \Vert f \Vert^*_{p,r} = \left( {r \over p}
   \int_0^\infty \big[ t^{1/p} f^*(t) \big]^r {dt \over t}
   \right)^{1/r} < \infty,
$$
where $1 \leq p < \infty$, $1 \leq r < \infty$. An operator $T$ is of
restricted weak $(p,p)$-type if $T :  L^{p,1}(\mu) \longrightarrow
L^{p,\infty}(\mu)$ is bounded, which is equivalent to $\Vert T\chi_E
\Vert_{L^p_*(\mu)} \leq C \Vert \chi_E \Vert_{L^p(\mu)}$ for all characteristic
functions $\chi_E$, with $C > 0$ independent of $E$. We refer the reader to
[11] for further information on these topics.

If both $\alpha, \beta > -1/2$, the authors proved (see [2]) that the $n$-th
partial sum operators are uniformly of restricted weak $(p,p)$-type but not of
of weak $(p,p)$-type when $p$ is an endpoint of the interval of mean
convergence. In theorems 2 and 3 we extend this result to weighted case $f
\longrightarrow u S_n(u^{-1} f)$, where $u$ is also a Jacobi weight, that is,
$$
   u(x) = (1-x)^a (1+x)^b, \qquad a, b \in \R.
$$
Now, the weighted uniform boundedness
$$
   \Vert uS_nf \Vert _{L^p(w)} \leq C \Vert uf \Vert _{L^p(w)}
   \quad \forall f \in L^p(u^pw), \ \forall n \geq 0
$$
holds (see [6]) if and only if
$$
   \eqalign{
   \vert a+(\alpha+1)({1 \over p} - {1 \over 2}) \vert
   < \min \{{1 \over 4}, {\alpha+1 \over 2}\},
   \cr
   \vert b+(\beta+1)({1 \over p} - {1 \over 2}) \vert
   < \min \{{1 \over 4}, {\beta+1 \over 2}\}.
   \cr} \eqno (2)
$$

Via Pollard's formula, these operators can be related to the Hilbert transform.
Then, the theory of $A_p$ weights is used, as well as some classical
dyadic-type decomposition of the interval $[-1,1]$.

In the general case $\alpha > -1/2$, $\alpha \geq \beta$ (the case $\beta \geq
\alpha$ follows by symmetry), we prove in theorem 1 that the $n$-th partial sum
operators are uniformly of restricted weak $(p,p)$-type when $p$ is an endpoint
of the interval of mean convergence, thus extending the above cited result (the
question of the weak boundedness had already been answered in the negative in
[2]). Now, however, uniform bounds are not available for Jacobi polynomials;
therefore, a uniform weighted norm inequality is needed for operators of the
form $f \longrightarrow u_n H(v_n f)$, where $H$ is the Hilbert transform and
$(u_n)$, $(v_n)$ are two sequences of weights involving Jacobi polynomials or
their bounds. This is achieved by studying the $A_p$ constants of the pairs of
weights $(u_n, v_n)$, as well as some $L^{p,\infty}$ norms.

Concerning mixed weak norm inequalities for the Hilbert transform, we can state
the following property, which can be proved in the same way as theorem 3 of
[7] (throughout this paper, the Hilbert transform, as well as $A_p$ classes of
weights, are taken on $[-1,1]$): assume that $u_1(x), u_2(x), v(x) \geq 0$,
$1<p<\infty$ and there is a constant $C>0$ such that
$$
   \Vert u_2 Hg \Vert_{L^p_*(u_1)} \leq C \Vert g \Vert_{L^p(v)}
   \quad \forall g \in L^p(v);
$$
then, there exists another constant $B>0$ which depends only on $C$, such that
for every interval $I$
$$
   \Vert u_2 \chi_I \Vert_{L^p_*(u_1)} \left(
   \int_{-1}^1 {v(x)^{-1/(p-1)} \over (\vert I \vert + \vert x-x_I \vert)^q} dx
   \right)^{1/q} \leq B, \eqno (3)
$$
$x_I$ being the centre of $I$.

Let us state the main results of this paper. By symmetry, there is no loss of
generality in assuming $\alpha \geq \beta$. Regarding the restricted weak
type, by standard arguments it is enough to consider just one of the endpoints
of the interval of mean convergence, as we remark below.

\proclaim{Theorem 1}. Let $\alpha > -1/2$, $\beta > -1$, $\alpha \geq \beta$. If
$p = {4(\alpha+1) \over 2\alpha+1}$, there exists a constant $C>0$ such that
for every measurable set $E$ and for every $n \geq 0$
$$
   \Vert S_n\chi_E \Vert _{L^p_*(w)} \leq	 C \Vert \chi_E \Vert _{L^p(w)}.
$$

\medskip

\proclaim{Theorem 2}. Let $\alpha, \beta \geq -1/2$, $u(x)=(1-x)^a(1+x)^b$,
$1<p<\infty$. If the inequalities
$$
   -{1 \over 4} \leq a+(\alpha+1)({1 \over p}-{1 \over 2}) < {1 \over 4},
   \qquad
   -{1 \over 4} \leq b+(\beta+1)({1 \over p}-{1 \over 2}) < {1 \over 4}
$$
hold, then there exists a constant $C>0$ such that
$$
   \Vert uS_n(u^{-1}\chi_E) \Vert _{L^p_*(w)}
   \leq C \Vert \chi_E \Vert _{L^p(w)}
$$
for every $n \geq 0$ and every measurable set $E \subseteq [-1,1]$.

\medskip

{\noindent\bf Remark.} For $1<p<\infty$ and $1/p+1/q=1$, it is easy to see that
$$
   \Vert uS_n(u^{-1}\chi_E) \Vert _{L^p_*(w)}
   \leq C \Vert \chi_E \Vert _{L^p(w)}
   \quad \forall n\geq 0, \ \forall E \subseteq [-1,1]
$$
if and only if
$$
   \Vert u^{-1}S_n(u\chi_E) \Vert _{L^q_*(w)}
   \leq C \Vert \chi_E \Vert _{L^q(w)}
   \quad \forall n\geq 0, \ \forall E \subseteq [-1,1].
$$
This allows us to derive, from theorem 2, the same result for the case
$$
   -{1 \over 4} < a+(\alpha+1)({1 \over p}-{1 \over 2}) \leq {1 \over 4},
   \qquad
   -{1 \over 4} < b+(\beta+1)({1 \over p}-{1 \over 2}) \leq {1 \over 4}
$$
as well as the analog of theorem 1 for $p={4(\alpha+1) \over 2\alpha+3}$.

\medskip

\proclaim{Theorem 3}. Let $\alpha, \beta \geq -1/2$, $u(x)=(1-x)^a(1+x)^b$,
$1<p<\infty$. If there exists a constant $C>0$ such that for every $f \in
L^p(u^pw)$ and for every $n \geq 0$
$$
   \Vert uS_nf \Vert _{L^p_*(w)} \leq	 C \Vert uf \Vert _{L^p(w)},
$$
then the inequalities
$$
   \vert a+(\alpha+1)({1 \over p}-{1 \over 2}) \vert < {1 \over 4}, \qquad
   \vert b+(\beta+1) ({1 \over p}-{1 \over 2}) \vert < {1 \over 4}
$$
are verified.

\bigskip

{\noindent \bf \S 2. Preliminary lemmas.}

\medskip

\noindent A basic tool in the study of Fourier series on the interval $[-1,1]$
is Pollard's decomposition of the kernels $K_n(x,t)$ (see [9], [6]): if
$\{P_n\}_{n\geq 0}$ is the sequence of polynomials orthonormal with respect to
$w(x)dx$ and $\{Q_n\}_{n\geq 0}$ is the sequence of polynomials associated to
\hbox{$(1-x^2)w(x)dx$}, then
$$
   K_n(x,t) = r_nT_{1,n}(x,t) + s_nT_{2,n}(x,t) + s_nT_{3,n}(x,t),
$$
where
$$
   T_{1,n}(x,t) = P_{n+1}(x) P_{n+1}(t),
$$
$$
   T_{2,n}(x,t) = (1-t^2) {P_{n+1}(x) Q_n(t) \over x-t},
$$
$$
   T_{3,n}(x,t) = (1-x^2) {P_{n+1}(t) Q_n(x) \over t-x}
$$
and $\{r_n\}$, $\{s_n\}$ are bounded sequences. In fact, for any measure $\mu$
on $[-1,1]$ with $\mu'>0$ $a.e.$ (in particular, for $w(x)dx$),
$$
   \vbox{\halign{\hfil #  &  \qquad \hfil # \hfil \cr
   $\displaystyle\lim_{n\to\infty} r_n = -1/2$,
   & $\displaystyle\lim_{n\to\infty} s_n = 1/2$ \cr
   }}
$$
(this can be deduced from [9] and [10] or [4]). Therefore, we can write
$$
   S_nf = r_n W_{1,n}f + s_n W_{2,n}f - s_n W_{3,n}f,
$$
where
$$
   W_{1,n}f(x) = P_{n+1}(x) \int_{-1}^1 P_{n+1}(t) f(t) w(t) dt,
$$
$$
   W_{2,n}f(x) = P_{n+1}(x) H((1-t^2) Q_n(t) f(t) w(t), x)
$$
and
$$
   W_{3,n}f(x) = (1-x^2) Q_n(x) H(P_{n+1}(t) f(t) w(t), x),
$$
$H$ being the Hilbert transform on the interval $[-1,1]$. Thus, the study of
$S_n$ can be reduced to that of $W_{i,n}$ $(i=1,2,3)$. Due to the definition of
$W_{2,n}$ and $W_{3,n}$, we will need to show the uniform boundedness of the
Hilbert transform with pairs of weights $(u_n,v_n)$.

The boundedness of the Hilbert transform can be stated in terms of Muckenhoupt's
classes of weights $A_p$ (see [3] and [8]; throughout this paper, these will be
$A_p$ classes on the interval $[-1,1]$): if $u$, $v$ are two weights,
$1<p<\infty$ and $(u^{\delta},v^{\delta}) \in A_p$ for some $\delta >1$, then
$H$ is a bounded operator from $L^p(v)$ into $L^p(u)$, with a constant which
depends only on the $A_p$ constant of $(u^{\delta},v^{\delta})$.

Therefore, we will say that a sequence $\{(u_n,v_n)\}_n$ belongs {\it uniformly}
to an $A_p$ class if $(u_n,v_n) \in A_p$ $\forall n$ with a constant that does
not depend on $n$.

The polynomials $P_n$ satisfy the estimate
$$
   \vert P_n(x) \vert \leq C (1-x+n^{-2})^{-(2\alpha+1)/4}
   (1+x+n^{-2})^{-(2\beta +1)/4} \quad \forall n, \forall x \in [-1,1] \eqno (4)
$$
with a constant $C>0$ independent of $x$ and $n$. An analogous estimate is
verified by $Q_n$, with $\alpha+1$ and $\beta+1$ instead of $\alpha$ and
$\beta$. In this case, as $2\alpha+3>0$ and $2\beta+3>0$, we can remove the
$n$'s and get
$$
   \vert Q_n(x) \vert \leq C (1-x)^{-(2\alpha+3)/4}
   (1+x)^{-(2\beta +3)/4} \quad \forall n, \forall x \in [-1,1]. \eqno (5)
$$
In this context, the following result will be useful.

\medskip

\proclaim{Lemma 4}. Let $\{x_n\}$ be a sequence of positive numbers with
$\displaystyle\lim_{n\to\infty} x_n = 0$. Let \hbox{$1<p<\infty$}, \hbox{$r,s,R,
S\in \R$}. Then,
$$
   (\vert x \vert^r (\vert x \vert + x_n)^s,
   \vert x \vert^R (\vert x \vert + x_n)^S) \in A_p([-1,1]) \ uniformly
$$
if and only if
$$
   \vbox {\halign{# \hfil  & \qquad \hfil #  &  \qquad \hfil # \hfil \cr
   $-1<r$; & $R < p-1$; & $R \leq r$; \cr
   $-1<r+s$; & $R+S < p-1$; & $R+S \leq r+s$. \cr
   }}
$$

\proof According to its definition,
$$
   (\vert x \vert^r (\vert x \vert + x_n)^s,
   \vert x \vert^R (\vert x \vert + x_n)^S) \in A_p([-1,1]) \ uniformly
$$
if and only if there exists a constant $C>0$ such that
$$
   \int_a^b \vert x \vert^r (\vert x \vert + x_n)^s dx
   \ (\int_a^b [\vert x \vert^R (\vert x \vert + x_n)^S]^{-1/(p-1)}dx)^{p-1}
   \leq C (b-a)^p \eqno (6)
$$
for all $-1\leq a < b \leq 1$, $\forall n \geq 1$. Integrability conditions
imply the above inequalities. In turn, if those inequalities hold we can easily
deduce (6) from the estimate
$$
   \int_a^b x^{\gamma} (x+c)^{\mu} dx \leq
   \cases{K b^{\gamma+\mu}(b-a) &if $c \leq b$ \cr
   K b^{\gamma}c^{\mu}(b-a) &if $b \leq c$},
$$
valid for $0 \leq a < b \leq 1$, $0 \leq c \leq 1$, $\gamma > -1$, $\gamma +
\mu > -1$, with a constant $K$ which depends only on $\gamma$, $\mu$.
\qed

\medskip

The same property holds if we replace $x$ by $x-a$, with $a \in [-1,1]$. Even
more, it is not difficult to show that in order to see whether a finite product
of this type of expressions is uniformly in $A_p$, we only need to check the
above inequalities for each factor of the form
$$
   (\vert x-a \vert^r (\vert x-a \vert + x_n)^s,
   \vert x-a \vert^R (\vert x-a \vert + x_n)^S)
$$
separately.

\medskip

We will eventually need to show that some of the operators are not of
strong or weak type. In this sense, the following lemma (see [5]) will be used:

\proclaim{Lemma 5}. Let ${\rm supp}\ d\alpha = [-1,1]$, $\alpha'>0$ a.e. in
$[-1,1]$, and $0<p \leq \infty$. There exists a constant $C>0$ such that if $g$
is a Lebesgue-measurable function on $[-1,1]$, then
$$
   \Vert \alpha'(x)^{-1/2}(1-x^2)^{-1/4} \Vert _{L^p(\vert g \vert^pdx)} \leq
   C \liminf_{n \to \infty} \Vert P_n \Vert _{L^p(\vert g \vert^pdx)}.
$$

\medskip

There is a weak version of this property: it is a consequence of Kolmogorov's
condition (see [1], lemma V.2.8, p. 485) and the previous lemma.

\proclaim{Lemma 6}. Let ${\rm supp}\ d\alpha = [-1,1]$, $\alpha'>0$ a.e. in
$[-1,1]$, and $0<p<\infty$. There exists a constant $C>0$ such that if
$g$, $h$ are Lebesgue-measurable functions on $[-1,1]$, then
$$
   \Vert \alpha'(x)^{-1/2}(1-x^2)^{-1/4}g(x) \Vert _{L^p_*(\vert h \vert^pdx)}
   \leq C \liminf_{n \to \infty} \Vert P_ng \Vert _{L^p_*(\vert h \vert^pdx)}.
$$

\medskip

The following lemma will be useful to estimate some weighted $L^p_*$ norms:

\proclaim{Lemma 7}. Let $1 \leq p <\infty$, $r,s \in \R$, $a>0$. Then,
$$
   \chi_{(0,a)}(x) x^r \in L^p_*(x^sdx) \Longleftrightarrow
   pr+s+1 \geq 0, \ \ (r,s) \not= (0,-1).
$$
Moreover, in this case there is a constant $K$ depending on $r,s,p$ such that
$$
   \Vert \chi_{(0,a)}(x) x^r \Vert _{L^p_*(x^sdx)} = K a^{r+(s+1)/p}.
$$

\proof Since
$$
   \Vert \chi_{(0,a)}(x) x^r \Vert ^p_{L^p_*(x^sdx)} =
   \sup_{y>0} y^p\int_A x^sdx
$$
with $A = \{x;\ 0<x<a,\ x^r>y\}$, the proof is reduced to a simple calculation
of that integral, depending on the sign of $pr+s+1$ and $r$. \qed

\medskip

Finally, this lemma will be used in the study of the operator $W_{2,n}$:

\proclaim{Lemma 8}. Let $\alpha >-1$, $1<p< \infty$, $1/p+1/q=1$, $0<r<1$, $n
\in \N$. If \hbox{$(\alpha +1)(1/p-1/2) < 1/4$}, then there exists a constant
$C$, independent of $r$ and $n$, such that
$$
   \Vert (1-t)(1-t+n^{-2})^{-(2\alpha+3)/4}\chi_{(r,1)}(t) \Vert
   _{L^q_*((1-t)^\alpha)}
   \leq C (1-r)^{1-(\alpha+1)/p}(1-r+n^{-2})^{(2\alpha+1)/4}.
$$

\proof a) Case $\alpha \geq -1/2$. Since $-(2\alpha+3)/4<0$,
$$
   \Vert (1-t)(1-t+n^{-2})^{-(2\alpha+3)/4}\chi_{(r,1)}(t) \Vert
   _{L^q_*((1-t)^\alpha)} \leq
$$
$$
   \leq \Vert (1-t)^{(1-2\alpha)/4}\chi_{(r,1)}(t) \Vert
   _{L^q_*((1-t)^\alpha)}.
$$
By lemma 7 and taking into account that $(2\alpha+1)/4 \geq 0$, we have
$$
   \Vert (1-t)^{(1-2\alpha)/4}\chi_{(r,1)}(t) \Vert
   _{L^q_*((1-t)^\alpha)} \leq C (1-r)^{(1-2\alpha)/4+(\alpha+1)/q} \leq
$$
$$
   \leq C (1-r)^{1-(\alpha+1)/p}(1-r+n^{-2})^{(2\alpha+1)/4}.
$$

b) Case $1-r \leq n^{-2}$. Then, by lemma 7 and the inequality $1-r+n^{-2} \leq
2n^{-2}$, we obtain
$$
   \Vert (1-t)(1-t+n^{-2})^{-(2\alpha+3)/4}\chi_{(r,1)}(t) \Vert
   _{L^q_*((1-t)^\alpha)} \leq
$$
$$
   \leq (n^{-2})^{-(2\alpha+3)/4} \Vert (1-t)\chi_{(r,1)}(t) \Vert
   _{L^q_*((1-t)^\alpha)} \leq
$$
$$
   \leq C (n^{-2})^{-(2\alpha+3)/4} (1-r)^{1+(\alpha+1)/q} \leq
$$
$$
   \leq C (1-r+n^{-2})^{-(2\alpha+3)/4} (1-r)^{1-(\alpha+1)/p} (1-r)^{\alpha+1}
   \leq 
$$
$$
   \leq C (1-r)^{1-(\alpha+1)/p}(1-r+n^{-2})^{(2\alpha+1)/4}.
$$

c) Case $\alpha < -1/2$ and $n^{-2} \leq 1-r$. Then $1-2\alpha \geq 0$ and
$1-r+n^{-2} \leq 2(1-r)$. Thus,
$$
   \Vert (1-t)(1-t+n^{-2})^{-(2\alpha+3)/4}\chi_{(r,1)}(t) \Vert
   _{L^q_*((1-t)^\alpha)} \leq
$$
$$
   \leq \Vert (1-t+n^{-2})^{(1-2\alpha)/4}\chi_{(r,1)}(t) \Vert
   _{L^q_*((1-t)^\alpha)} \leq
$$
$$
   \leq (1-r+n^{-2})^{(1-2\alpha)/4} \Vert \chi_{(r,1)}(t) \Vert
   _{L^q_*((1-t)^\alpha)} \leq
$$
$$
   \leq C (1-r+n^{-2})^{(1-2\alpha)/4}(1-r)^{(\alpha+1)/q} =
$$
$$
   = C (1-r)^{1-(\alpha+1)/p} (1-r+n^{-2})^{(2\alpha+1)/4}
   (1-r+n^{-2})^{-\alpha} (1-r)^{\alpha} \leq
$$
$$
   \leq C  (1-r)^{1-(\alpha+1)/p} (1-r+n^{-2})^{(2\alpha+1)/4}.\ \qed
$$

\bigskip

{\noindent \bf \S 3. Proof of theorems 1 and 2.}

\medskip

\noindent The proof of theorem 1 consists of lemmas 9, 10 and 11 below.
In order to prove theorem 2, analogous weighted lemmas can be shown using that,
in the case $\alpha, \beta \geq -1/2$, not only the polynomials $Q_n$ but also
the $P_n$ satisfy an uniform estimate similar to (5).

\proclaim {Lemma 9}. Under the hypothesis of theorem 1, there exists a constant
$C$ such that
$$
   \Vert W_{1,n}f \Vert _{L^p_*(w)} \leq C \Vert f \Vert _{L^p(w)}
   \qquad \forall f \in L^p(w), \ \forall n \in \N.
$$

\proclaim {Lemma 10}. Under the hypothesis of theorem 1, there exists a constant
$C$ such that
$$
   \Vert W_{3,n}f \Vert _{L^p(w)} \leq C \Vert f \Vert _{L^p(w)}
   \qquad \forall f \in L^p(w), \ \forall n \in \N.
$$

\proclaim {Lemma 11}. Under the hypothesis of theorem 1, there exists a constant
$C$ such that for every measurable set $E \subseteq [-1,1]$ and for every
$n \geq 0$
$$
   \Vert W_{2,n}\chi_E \Vert _{L^p_*(w)} \leq C \Vert \chi_E \Vert _{L^p(w)}
   \eqno (7)
$$

\medskip

{\it Proof of lemma 9}. From its definition, we have
$$
   \Vert W_{1,n}f \Vert _{L^p_*(w)} \leq \Vert P_{n+1} \Vert _{L^p_*(w)}
   \Vert P_{n+1} \Vert _{L^q(w)} \Vert f \Vert _{L^p(w)},
$$
where $1/p+1/q=1$. So, we only need to prove
$$
   \Vert P_n \Vert _{L^p_*(w)} \leq C \quad \forall n \in \N
$$
and
$$
   \Vert P_n \Vert _{L^q(w)} \leq C \quad \forall n \in \N,
$$
which follows from lemma 7, (4) and the dominate convergence
theorem. \qed

\medskip

{\it Proof of lemma 10}. It is clear that
$$
   \Vert W_{3,n}f \Vert _{L^p(w)} \leq C \Vert f \Vert _{L^p(w)} 
$$
for every $f \in L^p(w)$ if and only if
$$
   \Vert Hg \Vert _{L^p((1-x^2)^p \vert Q_n \vert^p w)} \leq
   C \Vert g \Vert _{L^p(\vert P_{n+1} \vert^{-p} w^{1-p})}
$$
for every $g \in L^p(\vert P_{n+1} \vert^{-p} w^{1-p})$. Using again (4) and
its analogous for $Q_n$, it is enough to obtain
$$
   \Vert Hg \Vert _{L^p(u_n)} \leq C \Vert g \Vert _{L^p(v_n)}
   \quad \forall n,\ \forall g \in L^p(v_n),
$$
with
$$
   u_n(x) = (1-x)^{p+\alpha}(1-x+n^{-2})^{-p(2\alpha+3)/4}
   (1+x)^{p+\beta}(1+x+n^{-2})^{-p(2\beta+3)/4}
$$
and
$$
   v_n(x) = (1-x)^{\alpha(1-p)}(1-x+n^{-2})^{p(2\alpha+1)/4}
   (1+x)^{\beta(1-p)}(1+x+n^{-2})^{p(2\beta+1)/4}.
$$
Now, we only need to prove that
$$
   ((1-x)^{(p+\alpha)\delta}(1-x+n^{-2})^{-p(2\alpha+3)\delta/4},
   (1-x)^{\alpha(1-p)\delta}(1-x+n^{-2})^{p(2\alpha+1)\delta/4})
   \in A_p
$$
and
$$
   ((1+x)^{(p+\beta)\delta}(1+x+n^{-2})^{-p(2\beta+3)\delta/4},
   (1+x)^{\beta(1-p)\delta}(1+x+n^{-2})^{p(2\beta+1)\delta/4})
   \in A_p
$$
uniformly in $n$, for some $\delta >1$. This can be deduced from lemma 4.
\qed

\medskip

{\it Proof of lemma 11}. From $p = {4(\alpha+1) \over 2\alpha+1}$ and $\alpha
\geq \beta$, it follows
$$
   -{1 \over 4} \leq (\alpha+1)({1 \over p} - {1 \over 2}) < {1 \over 4}
   \eqno(8)
$$
and
$$
   -{1 \over 4} \leq (\beta+1)({1 \over p} - {1 \over 2}) < {1 \over 4}.
   \eqno(9)
$$
We will prove that $(8)$ and $(9)$ imply $(7)$. By the symmetry of these
inequalities, we can consider only the case $E \subseteq [0,1]$.

a) We will show first that there exists a constant $C>0$ such that 
$$
   \Vert \chi_{[-3/4,3/4]}W_{2,n}(\chi_E) \Vert _{L^p_*(w)} \leq
   C \Vert \chi_E \Vert _{L^p(w)}
$$
for all $n \in \N$ and every measurable set $E \subseteq [0,1]$.

Since $1-x$ and $1+x$ are bounded away from $0$ and $\infty$ on $[-3/4,3/4]$,
from (4) and the definition of $W_{2,n}$ we get
$$
   \Vert \chi_{[-3/4,3/4]}W_{2,n}(\chi_E) \Vert _{L^p_*(w)} \leq
   C \Vert \chi_{[-3/4,3/4]} H((1-t^2)\chi_EQ_nw) \Vert
   _{L^p((1-x)^r(1+x)^s)}
$$
for any previously fixed $r$, $s$. If we find $r$, $s$ such that
$$
   H:L^p((1-x)^{\alpha-p(2\alpha+1)/4}(1+x)^{\beta-p(2\beta+1)/4})
   \longrightarrow L^p((1-x)^r(1+x)^s) \eqno (10)
$$
is bounded, then from (5) it would follow
$$
   \Vert \chi_{[-3/4,3/4]} H((1-t^2)\chi_EQ_nw) \Vert _{L^p((1-x)^r(1+x)^s)}
   \leq \Vert H((1-t^2)\chi_EQ_nw) \Vert _{L^p((1-x)^r(1+x)^s)} \leq
$$
$$
   \leq C \Vert (1-x^2)\chi_EQ_nw \Vert 
   _{L^p((1-x)^{\alpha-p(2\alpha+1)/4}(1+x)^{\beta-p(2\beta+1)/4})} \leq
   C \Vert \chi_E \Vert _{L^p(w)},
$$
as we want to show. In order to get $(10)$, it is enough to have
$$
   ((1-x)^{r\delta}(1+x)^{s\delta},
   (1-x)^{\delta[\alpha-p(2\alpha+1)/4]}(1+x)^{\delta[\beta-p(2\beta+1)/4]})
   \in A_p
$$
for some $\delta >1$. This is equivalent, by lemma 4, to the following
conditions: $$
   \vbox{\halign{# \hfil & \qquad # \hfil \cr
   $-1<r$; & $-1<s$; \cr
   $\alpha-p(2\alpha+1)/4<p-1$; & $\beta-p(2\beta+1)/4<p-1$; \cr
   $\alpha-p(2\alpha+1)/4 \leq r$; & $\beta-p(2\beta+1)/4 \leq s$. \cr
   }}
$$
It is easy to see that the second row inequalities hold, while for the
others we only need to take $r$ and $s$ large enough.

b) Now, we are going to prove that there exists a constant $C>0$ such that
$$
   \Vert \chi_{[-1,-3/4]}W_{2,n}(\chi_E) \Vert _{L^p_*(w)} \leq
   C \Vert \chi_E \Vert _{L^p(w)}
$$
for all $n \in \N$ and every measurable set $E \subseteq [0,1]$.

As $E \subseteq [0,1]$, we can drop the denominator $x-t$ in
$W_{2,n}(\chi_E)$ and, using the inequality $(5)$, we get
$$
   \vert \chi_{[-1,-3/4]}(x) W_{2,n}(\chi_E,x) \vert \leq
$$
$$
   \leq C \chi_{[-1,-3/4]}(x) \vert P_{n+1}(x) \vert
   \int_{-1}^1 (1-t)^{1-(2\alpha+3)/4} \chi_E(t) w(t) dt \leq
$$
$$
   \leq C \chi_{[-1,-3/4]}(x) \vert P_{n+1}(x) \vert
   \ \Vert (1-t)^{(1-2\alpha)/4} \Vert_{L^q(w)}
   \Vert \chi_E \Vert_{L^p(w)} \leq
$$
$$
   \leq C \Vert \chi_E \Vert_{L^p(w)}
   \chi_{[-1,-3/4]}(x) \vert P_{n+1}(x) \vert.
$$
Therefore
$$
   \Vert \chi_{[-1,-3/4]}W_{2,n}(\chi_E) \Vert_{L^p_*(w)} \leq
$$
$$
   \leq C \Vert \chi_E \Vert_{L^p(w)}
   \ \Vert \chi_{[-1,-3/4]} (1+x+n^{-2})^{-(2\beta+1)/4} \Vert_{L^p_*(w)}
   \leq C \Vert \chi_E \Vert_{L^p(w)},
$$
by the dominate convergence and lemma 7.

c) We must show now that there exists a constant $C>0$ such that
$$
   \Vert \chi_{[3/4,1]}W_{2,n}(\chi_E) \Vert _{L^p_*(w)} \leq
   C \Vert \chi_E \Vert _{L^p(w)}
$$
for all $n \in \N$ and every measurable set $E \subseteq [0,1]$. Let us define,
for $k=2,3,\dots$ the sets
$$
   I_k = [1-2^{-k},1-2^{-k-1}),
$$
$$
   J_{k1} = [0,1-2^{-k+1}), \qquad J_{k2} = [1-2^{-k+1},1-2^{-k-2}), \qquad
   J_{k3} = [1-2^{-k-2},1].
$$
For each $k\geq 2$, $J_{ki}$ $(i=1,2,3)$ are non-overlapping sets such that
$[0,1] = J_{k1} \cup J_{k2} \cup J_{k3}$. The sets $I_k$ are also
disjoint and $\displaystyle\bigcup_{k\geq 2} I_k = [3/4,1)$. The
following properties are easy to check:
$$
   \forall k\geq 2,\ \forall x\in I_k,\quad 2^{-k-1} \leq 1-x \leq 2^{-k};
   \eqno (11)
$$
$$
   \forall k\geq 2,\ \forall t\in J_{k2},
   \quad 2^{-k-2} \leq 1-t \leq 2^{-k+1}; \eqno (12)
$$
$$
   \forall k\geq 2,\ \forall x\in I_k,\ \forall t\in J_{k1},
   \quad 2^{-k+1} \leq 1-t \leq 2(x-t) \leq 2(1-t); \eqno (13)
$$
$$
   \forall k\geq 2,\ \forall x\in I_k,\ \forall t\in J_{k3},
   \quad 1-t \leq 2^{-k-2} \leq t-x \leq 2^{-k}. \eqno (14)
$$
We can write
$$
   \chi_{[3/4,1]} W_{2,n}(\chi_E) =
   \sum_{k\geq 2} \chi_{I_k}W_{2,n}(\chi_E\chi_{J_{k1}})
   + \sum_{k\geq 2} \chi_{I_k}W_{2,n}(\chi_E\chi_{J_{k2}})
   + \sum_{k\geq 2} \chi_{I_k}W_{2,n}(\chi_E\chi_{J_{k3}}).
$$
We prove that each term is bounded.

c1) If $x \in I_k$, from $(13)$ and H\"older's inequality for Lorentz spaces it
follows
$$
   \vert H((1-t^2)\chi_E \chi_{J_{k1}}Q_nw,x) \vert
   \leq C \int_{-1}^1 \chi_E \chi_{J_{k1}}\vert Q_n \vert w \leq
$$
$$
   \leq C \Vert \chi_E \Vert _{L^p(w)}
   \Vert \chi_{J_{k1}} Q_n \Vert _{L^q_*(w_)}.
$$
From the estimates $(4)$ for $Q_n$, property $(8)$, lemma 7 and using that $1
\leq 1+t \leq 2$ for $t \in J_{k1}$, we obtain
$$
   \Vert \chi_{J_{k1}} Q_n \Vert _{L^q_*(w_)} \leq C
   \Vert \chi_{J_{k1}}
   (1-t+n^{-2})^{(\alpha+1)/q-(2\alpha+3)/4}(1-t+n^{-2})^{-(\alpha+1)/q}
   \Vert _{L^q_*(w_)} \leq
$$
$$
   \leq C (1-x+n^{-2})^{(\alpha+1)/q-(2\alpha+3)/4}
   \Vert \chi_{J_{k1}} (1-t+n^{-2})^{-(\alpha+1)/q} \Vert _{L^q_*(w_)} \leq
$$
$$
   \leq C (1-x+n^{-2})^{(\alpha+1)/q-(2\alpha+3)/4}.
$$
Therefore, if $x \in I_k$ then
$$
   \vert H((1-t^2)\chi_E \chi_{J_{k1}}Q_nw,x) \vert
   \leq C \Vert \chi_E \Vert _{L^p(w)}
   (1-x+n^{-2})^{(\alpha+1)/q-(2\alpha+3)/4},
$$
where the constant $C$ does not depend on $k$. Since the $I_k$ are disjoint,
this implies $$
   \vert \sum_{k\geq 2} \chi_{I_k}(x) W_{2,n}(\chi_E\chi_{J_{k1}},x) \vert
   \leq
$$
$$
   \leq C (1-x+n^{-2})^{-(2\alpha+1)/4} \Vert \chi_E \Vert _{L^p(w)}
   (1-x+n^{-2})^{(\alpha+1)/q-(2\alpha+3)/4} =
$$
$$
   = C (1-x+n^{-2})^{-(\alpha+1)/p} \Vert \chi_E \Vert _{L^p(w)}
   \leq C (1-x)^{-(\alpha+1)/p} \Vert \chi_E \Vert _{L^p(w)}.
$$
Then, by lemma 7,
$$
   \Vert \sum_{k\geq 2} \chi_{I_k} W_{2,n}(\chi_E\chi_{J_{k1}}) \Vert
   _{L^p_*(w)} \leq C \Vert \chi_E \Vert _{L^p(w)}.
$$

c2) Let $k\geq 2$. By $(4)$ and $(11)$,
$$
   \Vert \chi_{I_k}W_{2,n}(\chi_E\chi_{J_{k2}}) \Vert_{L^p_*(w)} =
   C \Vert \chi_{I_k}P_{n+1}H((1-t^2)\chi_E\chi_{J_{k2}}Q_nw)
   \Vert_{L^p_*(w)} \leq
$$
$$
   \leq C (2^{-k}+n^{-2})^{-(2\alpha+1)/4}(2^{-k})^{\alpha/p} \Vert
   \chi_{I_k}H((1-t^2)\chi_E\chi_{J_{k2}}Q_nw) \Vert_{L^p_*(dx)} \leq
$$
$$
   \leq C (2^{-k}+n^{-2})^{-(2\alpha+1)/4}(2^{-k})^{\alpha/p}
   \Vert H((1-t^2)\chi_E\chi_{J_{k2}}Q_nw) \Vert_{L^p(dx)}.
$$
Since the Hilbert transform is bounded in $L^p(dx)$, this expression can be
bounded, using $(4)$ and $(12)$, by
$$
   C (2^{-k}+n^{-2})^{-(2\alpha+1)/4}(2^{-k})^{\alpha/p}
   \Vert (1-x^2)\chi_E\chi_{J_{k2}}Q_nw \Vert_{L^p(dx)} \leq
$$
$$
   \leq C
   (2^{-k}+n^{-2})^{-(2\alpha+1)/4}(2^{-k})^{\alpha+1}
   (2^{-k}+n^{-2})^{-(2\alpha+3)/4}
   \Vert (1-x)^{\alpha/p}\chi_E\chi_{J_{k2}} \Vert_{L^p(dx)} \leq
$$
$$
   \leq C \Vert \chi_E\chi_{J_{k2}} \Vert_{L^p(w)}.
$$
Now, as the functions $\chi_{I_k}W_{2,n}(\chi_E\chi_{J_{k2}})$ have
non-overlapping support and $\sum_{k\geq 2} \chi_{J_{k2}} \leq 3$, we get
$$
   \Vert \sum_{k\geq 2} \chi_{I_k}W_{2,n}(\chi_E\chi_{J_{k2}})
   \Vert ^p_{L^p_*(w)} \leq
   \sum_{k\geq 2} \Vert \chi_{I_k}W_{2,n}(\chi_E\chi_{J_{k2}})
   \Vert ^p_{L^p_*(w)} \leq
$$
$$
   \leq C \sum_{k\geq 2} \Vert \chi_E\chi_{J_{k2}} \Vert ^p_{L^p(w)}
   \leq C \Vert \chi_E \Vert ^p_{L^p(w)}.
$$
That is,
$$
   \Vert \sum_{k\geq 2} \chi_{I_k}W_{2,n}(\chi_E\chi_{J_{k2}})
   \Vert _{L^p_*(w)} \leq C \Vert \chi_E \Vert _{L^p(w)}.
$$

c3) Let $k\geq 2$ and $x \in I_k$. By $(14)$, H\"older's inequality for Lorentz
spaces and $(4)$, it follows
$$
   \vert H((1-t^2)\chi_E\chi_{J_{k3}}Q_nw,x) \vert \leq
$$
$$
   \leq C 2^k \int_{-1}^1(1-t^2)\chi_E(t)\chi_{J_{k3}}(t)Q_n(t)w(t)dt \leq
$$
$$
   \leq C 2^k \Vert (1-t)\chi_E\chi_{J_{k3}}Q_n \Vert _{L^1(w)} \leq
$$
$$
   \leq C 2^k \Vert \chi_E \Vert _{L^p(w)} \Vert
   (1-t)(1-t+n^{-2})^{-(2\alpha+3)/4} \chi_{J_{k3}}
   \Vert _{L^q_*((1-t)^{\alpha})}.
$$
By lemma 8 and (8),
$$
   \Vert (1-t)(1-t+n^{-2})^{-(2\alpha+3)/4} \chi_{J_{k3}} \Vert
   _{L^q_*((1-t)^{\alpha})} \leq C (2^{-k})^{1-(\alpha+1)/p}
   (2^{-k}+n^{-2})^{(2\alpha+1)/4},
$$
what, together with (11), implies
$$
   \vert H((1-t^2)\chi_E\chi_{J_{k3}}Q_nw,x) \vert \leq
   C (1-x)^{-(\alpha+1)/p} (1-x+n^{-2})^{(2\alpha+1)/4}
   \Vert \chi_E \Vert _{L^p(w)}
$$
if $x \in I_k$, with a constant $C$ which does not depend on $x$, $E$, $k$,
$n$. Since the $I_k$ are non-overlapping, we have
$$
   \vert \sum_{k\geq 2} \chi_{I_k}(x) W_{2,n}(\chi_E\chi_{J_{k3}},x) \vert
   \leq C \chi_{[3/4,1)} (1-x)^{-(\alpha+1)/p} \Vert \chi_E \Vert _{L^p(w)}
$$
and, by lemma 7,
$$
   \Vert \sum_{k\geq 2} \chi_{I_k} W_{2,n}(\chi_E\chi_{J_{k3}}) \Vert
   _{L^p_*(w)} \leq C \Vert \chi_E \Vert _{L^p(w)}.
$$
This concludes the proof of the lemma. \qed

\bigskip

{\noindent \bf \S 4. Proof of theorem 3.}

\medskip

\noindent The weak boundedness
$$
   \Vert uS_nf \Vert _{L^p_*(w)} \leq	 C \Vert uf \Vert _{L^p(w)}
$$
implies the following conditions (see [2], theorem 1, with the appropriate
changes):
$$
   u \in L^p_*(w)
$$
$$
   u^{-1} \in L^q(w)
$$
$$
   u(x) w(x)^{-1/2}(1-x^2)^{-1/4} \in L^p_*(w)
$$
$$
   u(x)^{-1} w(x)^{-1/2}(1-x^2)^{-1/4} \in L^q(w).
$$
With the weight $u(x) = (1-x)^a(1+x)^b$ and having in mind that $\alpha, \beta
\geq -1/2$, this means
$$
   \eqalign{
   -{1 \over 4} \leq a+(\alpha+1)({1 \over p}-{1 \over 2}) < {1 \over 4}, \cr
   -{1 \over 4} \leq b+(\beta+1)({1 \over p}-{1 \over 2}) < {1 \over 4}.
   \cr}
$$
Therefore, we only need to show that the equality cannot occur in the left hand
side of these equations. Assume, for example,
$$
   -{1 \over 4} = a+(\alpha+1)({1 \over p}-{1 \over 2}). \eqno (15)
$$
Let us consider again Pollard's decomposition of the partial sums $S_nf$. As we
mentioned at the beginning of the previous section, the proofs of lemmas 9 and
10 essentially show that there exists a constant $C$ such that for all $f \in
L^p(w)$ and every $n \in \N$
$$
   \Vert uW_{1,n}f \Vert _{L^p_*(w)} \leq	 C \Vert uf \Vert _{L^p(w)}
$$
and
$$
   \Vert uW_{3,n}f \Vert _{L^p(w)} \leq	 C \Vert uf \Vert _{L^p(w)}
$$
(notice that in the case $\alpha, \beta \geq -1/2$, the polynomials $P_n$
satisfy the estimate analogous to (5), what simplifies the proofs). Under our
hypothesis, this implies that there exists also a constant $C$ such that for all
$f \in L^p(u^pw)$ and every $n \in \N$
$$
   \Vert uW_{2,n}f \Vert _{L^p_*(w)} \leq	 C \Vert uf \Vert _{L^p(w)},
$$
that is,
$$
   \Vert u P_{n+1} Hg \Vert _{L^p_*(w)} \leq	 C
   \Vert u(x) (1-x^2)^{-1} Q_n(x)^{-1} w(x)^{-1} g \Vert _{L^p(w)}.
$$
Applying (3), we have
$$
   \Vert u P_{n+1} \chi_I \Vert _{L^p_*(w)} \left(
   \int_{-1}^1 { u(x)^{-q} (1-x^2)^q \vert Q_n(x) \vert^q w(x) \over
   (\vert I \vert + \vert x-x_I \vert)^q } dx
   \right)^{1/q} \leq C
$$
for every interval $I \subseteq [-1,1]$, with a constant $C>0$ independent of
$n$ and $I$; now, by lemma 6
$$
   \Vert u(x) (1-x^2)^{-1/4} w(x)^{-1/2} \chi_I \Vert _{L^p_*(w)} \left(
   \int_{-1}^1 { u(x)^{-q} (1-x^2)^{q/4} w(x)^{1-q/2} \over
   (\vert I \vert + \vert x-x_I \vert)^q } dx
   \right)^{1/q} \leq C.
$$
Taking $I = [1-\varepsilon,1]$, it follows
$$
   \Vert x^{a-\alpha/2-1/4} \chi_{[0,\varepsilon]} \Vert _{L^p_*(x^\alpha)}
   \left( \int_0^1 { x^{-aq+q/4+\alpha(1-q/2)} \over
   (\varepsilon + \vert x-\varepsilon/2 \vert)^q } dx
   \right)^{1/q} \leq C. \eqno (16)
$$
Now, by lemma 7 and (15)
$$
   \Vert x^{a-\alpha/2-1/4} \chi_{[0,\varepsilon]} \Vert _{L^p_*(x^\alpha)}
   = K \eqno (17)
$$
and
$$
   \int_0^1 { x^{-aq+q/4+\alpha(1-q/2)} \over
   (\varepsilon + \vert x-\varepsilon/2 \vert)^q } dx =
   \int_0^1
   { x^{1/(p-1)} \over (\varepsilon + \vert x-\varepsilon/2 \vert)^q } dx \geq
   C \int_\varepsilon^1 x^{1/(p-1) - q} dx = C \vert \log \varepsilon \vert,
$$
which, together with (17), leads to a contradiction in (16). Therefore, (15)
cannot be true and the theorem is proved.

\bigskip

\centerline {\bf REFERENCES}

\bigskip

\frenchspacing

\item {\bf [1]} {\cmcsc J. Garc\'\i a Cuerva} and {\cmcsc J. L. Rubio de
Francia}, {\it \lq \lq Weighted norm inequalities and related topics"},
North-Holland, Amsterdam, 1985.

\item {{\bf [2]}} {\cmcsc J. J. Guadalupe, M. P\'erez} and {\cmcsc J. L.
Varona}, {\it Weak behaviour of Fourier-Jacobi series}, J. Approx. Theory
{\bf 61} (1990), 222--238.

\item {{\bf [3]}} {\cmcsc R. Hunt, B. Muckenhoupt} and {\cmcsc R. Wheeden}, {\it
Weighted norm inequalities for the conjugate function and Hilbert transform},
Trans. Amer. Math. Soc. {\bf 176} (1973), 227--251.

\item {{\bf [4]}} {\cmcsc A. M\'at\'e, P. Nevai} and {\cmcsc V. Totik}, {\it
Asymptotics for the ratio of the leading coefficients of orthogonal polynomials
on the unit circle}, Constr. Approx. {\bf 1} (1985), 63--69.

\item {{\bf [5]}} {\cmcsc A. M\'at\'e, P. Nevai} and {\cmcsc V. Totik}, {\it
Necessary conditions for weighted mean convergence of Fourier series in
orthogonal polynomials}, J. Approx. Theory {\bf 46} (1986), 314--322.

\item {{\bf [6]}} {\cmcsc B. Muckenhoupt}, {\it Mean convergence of Jacobi
series}, Proc. Amer. Math. Soc. {\bf 23} (1969), 306--310.

\item {{\bf [7]}} {\cmcsc B. Muckenhoupt} and {\cmcsc R. L. Wheeden}, {\it Two
weight function norm inequalities for the Hardy-Littlewood maximal function and
the Hilbert transform}, Studia Math. {\bf 55} (1976), 279--294.

\item {{\bf [8]}} {\cmcsc C. J. Neugebauer}, {\it Inserting $A_p$ weights},
Proc. Amer. Math. Soc. {\bf 87} (1983), 644--648.

\item {{\bf [9]}} {\cmcsc H. Pollard}, {\it The mean convergence of orthogonal
series. II}, Trans. Amer. Math. Soc. {\bf 63} (1948), 355--367.

\item {{\bf [10]}} {\cmcsc E. A. Rahmanov}, {\it On the asymptotics of the ratio
of orthogonal polynomials. II}, Math. USSR. Sb. {\bf 46} (1983), 105--117.

\item {\bf [11]} {\cmcsc E. M. Stein} and {\cmcsc G. Weiss}, {\it \lq \lq
Introduction to Fourier analysis on euclidean spaces"}, Princeton Univ. Press,
New Jersey, 1975.


\bye
